% Original PDFToolkit regression fixture, distributed under the repository MIT licence.
% Compile twice with LuaLaTeX from TeX Live 2026. No shell escape is required.
% Fixed identifiers and random seed make the retained fixture byte reproducible.
\ExplSyntaxOn\sys_gset_rand_seed:n {20261029}\ExplSyntaxOff
\pdfvariable trailerid {[<504446546f6f6c6b69744d617468534501> <504446546f6f6c6b69744d617468534501>]}
\DocumentMetadata{
  lang=en-US,
  pdfversion=2.0,
  tagging=on,
  tagging-setup={math/setup=mathml-SE}
}
\documentclass{article}
\usepackage{unicode-math}
\pagestyle{empty}
\title{Mathematics in the PDF structure tree}
\author{PDFToolkit}
\date{}
\begin{document}
\maketitle
\thispagestyle{empty}

\section{Inline mathematics}
The expression $x^{2}+\frac{1}{2}$ stays inside this sentence.

\section{Displayed mathematics}
The fraction and exponent below are also mathematical structure elements.
\[
E=mc^{2}+\frac{1}{2}
\]

\section{A small matrix}
Rows and columns are represented by MathML table elements.
\[
\begin{pmatrix}1&2\\3&4\end{pmatrix}
\]

The source PDF contains MathML tags, without an associated MathML XML file.
\end{document}
